\vssub
\subsubsection{~曲线网格} \label{sub:num_space_curv}
\conthead{\ws (NRL Stennis)}{W. E. Rogers, T. J. Campbell}

\noindent 
作为传统“规则”网格的扩展，\ws\ 可在“不规则”网格上进行计算。这使模式可运行于替代网格投影（例如 Lambert 正形圆锥投影）、旋转网格，或近岸分辨率更高的贴岸线网格；不过，条件稳定格式对时间步长的限制仍然适用。不规则网格使用的传播格式与规则网格相同（\para\ref{sub:num_space_trad}）。

其实现由 \cite{rep:RC09} 详细描述，这里概述如下：使用 Jacobian 在整个区域内将正常的曲线空间与拉直空间相互转换。该转换只在传播例程内部进行，而不是在拉直空间中积分整个模式。每次调用传播子程序时（即每个时间步和每个谱分量），采用一个简单的三步过程：第一，用 Jacobian 把因变量（波作用量密度）转换到拉直空间；第二，通过对两个网格轴分别调用子程序来传播波作用量密度；第三，把波作用量密度转换回正常的曲线空间。实际通量计算相对于原来的规则网格形式没有显著改变。无论网格类型（规则或不规则）如何，都采用同一过程；对于规则网格，Jacobian 为 1。

关于用户接口：在 {\file ww3\_grid.inp} 中，用字符串表示网格类型。当该网格字符串为 `{\code RECT}' 时，模式处理规则网格输入。当该网格字符串为 `{\code CURV}' 时，模式处理不规则网格输入。[注意，在 \ws\ 4.00 版中，坐标系统（即度或米）和闭合类型（例如全球/环绕网格）也在 {\file ww3\_grid.inp} 中指定；开关 {\code LLG} 和 {\code XYG} 已废弃。]

从 \ws\ 第 5 版起，增加了在一种特殊曲线网格即“三极网格”上使用一阶传播格式运行的能力。在北半球，该网格类型使用两个极点而不是一个，并且两个极点都位于陆地上，以避免网格距奇异性。这类网格有时用于海洋模式，例如 \citep{art:Murr96} 和 \citep{art:Metz14}。不需要特殊开关，网格像其他不规则网格一样读入，但用户必须在 {\file ww3\_grid.inp} 中指定闭合类型 ({\code CSTRG}) 为 {\code TRPL}。具体细节见 \para\ref{sub:ww3grid} 中 {\file ww3\_grid.inp} 的文档。传播和梯度计算已修改，以处理新的闭合方法。{\code TRPL} 闭合类型仅与一阶 {\code PR1} 传播格式兼容。三极网格的一个吸引人特性是，它允许用户运行一个一直延伸到北极的单一网格。不过，虽然三个极点都位于陆地上，最靠近它们的海点处仍存在经线汇聚，这意味着以真实距离表示的网格距（决定最大传播时间步长）仍高度变化。通过使用多网格能力，可以获得更高效的网格距（即真实距离上的网格距变化较小）。虽然该格式处理了极点处网格距的奇异性，但并未处理与波向定义相关的奇异性。
